\section{Theory}\label{app:theory}

\subsection{Deriving the Classical Light-Matter Hamiltonian}\label{app:hamiltonian}
% [Muhammad: Done - Added lead-in; derivation follows Tao/Li et al., full derivation kept]
In this section, we present the light-matter Hamiltonian following the approach of Li, Subotnik, and Nitzan.\cite{li2020cavity,li2021cavity} 
The derivation below is included for completeness, with particular emphasis on obtaining a classical treatment suitable for molecular dynamics (MD) simulations. 
By treating the system classically and applying suitable approximations, one can study strong light-matter coupling for many molecules, a crucial requirement for glassy dynamics.
We begin with the minimal coupling Hamiltonian $\hat{H}$ describing the interaction between matter, composed of negatively charged electrons and positively charged nuclei, and an electromagnetic (EM) field:
\begin{gather}
\hat{H} = \hat{H}_\mathrm{M}+\hat{H}_\mathrm{EM} \,, \label{eq:H}
\\
\hat{H}_\mathrm{M} = \sum_{i=1}^{N_\mathrm{n}} \frac{| \hat{\bm P}_i -Z_i e \hat{\bm A}(\bm R_i)|^2}{2 M_i} + \sum_{j=1}^{N_e} \frac{| \hat{\bm p}_j +e \hat{\bm A}(\bm r_j)|^2}{2 m_j} + \hat{V}_\mathrm{nn}+ \hat{V}_\mathrm{ee}+ \hat{V}_\mathrm{ne} \,, \label{eq:HM}
\\
\hat{H}_\mathrm{EM} = \frac{\epsilon_0}{2}\int_\Omega d^3 \bm r  \ \big( | \hat{\bm E}(\bm r)|^2+c^2 | \hat{\bm  B}(\bm r)|^2 \big) \,, \label{eq:HF}
\end{gather}
Here, $\hat{H}_\mathrm{EM}$ and $\hat{H}_\mathrm{M}$ represent the electromagnetic field and matter components of the Hamiltonian, respectively.
The first two terms in $\hat{H}_\mathrm{M}$ describe the kinetic energy of $N_\mathrm{n}$ nuclei (charge $+Z_i e$, mass $M_i$) and $N_e$ electrons (charge $-e$, mass $m_j$) coupled to the vector potential $\hat{\bm A}$ of the EM field.
The remaining terms $\hat{V}_\mathrm{nn}$, $\hat{V}_\mathrm{ee}$, and $\hat{V}_\mathrm{ne}$ represent nucleus-nucleus, electron-electron, and nucleus-electron Coulombic interactions, respectively, generated by the scalar potential $\phi$ of the electromagnetic field.
In writing \cref{eq:H,eq:HF}, we have implicitly adopted the Coulomb gauge where $\nabla \cdot \hat{\bm A} = 0$, and thus the EM fields in \cref{eq:HF} are purely transverse.

We transform the field Hamiltonian into a harmonic oscillator form. 
To this end, we write the EM energy density in terms of the vector potential, using the relations $\hat{\bm E} = -\partial_t \hat{\bm A}$ and $\hat{\bm B} = \nabla \times \hat{\bm A}$ to obtain:
\begin{equation}
\hat{H}_\mathrm{EM}
= \frac{\epsilon_0}{2} \int_\Omega d^3 \bm r \,\Big(\bigl|\partial_t \hat{\bm A} \bigr|^2
+ c^2 \bigl|\nabla \times \hat{\bm A}\bigr|^2 \Big).
\label{eq:Hem_inA}
\end{equation}
Next, we consider a one-dimensional Fabry--Pérot cavity with mirrors at $z=\pm L/2$ and expand the transverse vector potential in terms of normal mode functions. We set the normal amplitudes as $\hat{C}_{\ell,\alpha}$ and write the expansion as
\begin{equation}
\hat{\bm A}(\bm r,t)
= \sum_{\ell,\alpha} \bigl(\hat{C}_{\ell,\alpha} (t)
+ \hat{C}_{\ell,\alpha}^\dagger \bigr)\,
\hat{\bm e}_\alpha\, u_\ell(z),
\end{equation}
where $\hat{C}_{\ell,\alpha}(t) \sim e^{-\mathrm{i}\omega_\ell t}$ and $\hat{\bm e}_\alpha \perp \hat{\bm z}$ are fixed polarization vectors and 
\begin{equation}
u_\ell(z) = \sqrt{\frac{2}{\Omega}}\,\sin\!\left[k_\ell\!\left(z+\frac{L}{2}\right)\right], \qquad
k_\ell=\frac{\ell\pi}{L}, \qquad \Omega = AL.
\end{equation}
The mode functions satisfy orthogonality,
\begin{equation}
\int_\Omega d^3\bm r\,(\hat{\bm e}_\alpha \!\cdot\! \hat{\bm e}_{\alpha^\prime})\,u_\ell(z)\,u_{\ell'}(z)
= \delta_{\ell\ell'}\,\delta_{\alpha\alpha'} \,.
\end{equation}
% An important identity for later use is the derivative of the mode functions, which also satisfy orthogonality %and for the sines with nodes at the mirrors one also has
% \begin{equation}
% \int_\Omega d^3\bm r\, u_\ell'(z)\,u_{\ell'}'(z) = k_\ell^2\,\delta_{\ell\ell'}.
% \end{equation}
Now, we can write the time derivative and curl of the vector potential as 
\begin{gather}
\partial_t \hat{\bm A}(\bm r,t)
= \sum_{\ell,\alpha} \mathrm{i}\omega_\ell\bigl(\hat{C}_{\ell,\alpha}^\dagger(t)-\hat{C}_{\ell,\alpha}(t)\bigr)\,
\hat{\bm e}_\alpha\,u_\ell(z), \label{eq:partialtA}\\
\nabla\times \hat{\bm A}(\bm r,t)
= \sum_{\ell,\alpha} \bigl(\hat{C}_{\ell,\alpha}(t)+\hat{C}_{\ell,\alpha}^\dagger(t)\bigr)\,
\hat{\bm e}_\alpha'\,u_\ell'(z), \label{eq:curlA}
\end{gather}
where $\hat{\bm e}_\alpha' \equiv \hat{\bm z}\times \hat{\bm e}_\alpha$ is $\hat{\bm e}_\alpha$ rotated by $90^\circ$. Substituting \cref{eq:partialtA,eq:curlA} into \cref{eq:Hem_inA} and using the orthogonality relations together with $\omega_\ell = c k_\ell$ yield
\begin{equation}
\hat{H}_\mathrm{EM}
= \epsilon_0 \sum_{\ell,\alpha}\omega_\ell^2
\left[ \hat{C}_{\ell,\alpha}^\dagger \hat{C}_{\ell,\alpha}+\hat{C}_{\ell,\alpha} \hat{C}_{\ell,\alpha}^\dagger\right].
\end{equation}
where we have suppressed the time-dependence for brevity. As is customary,  each cavity mode is associated with annihilation/creation operators $\hat{a}_{\ell,\alpha}$, $\hat{a}_{\ell,\alpha}^\dagger$ via
\begin{equation}
\hat{C}_{\ell,\alpha}=\sqrt{\frac{\hbar}{2\epsilon_0\omega_\ell}}\ \hat{a}_{\ell,\alpha} \,.
\end{equation}
With this choice, the Hamiltonian turns into
\begin{equation}
\hat{H}_\mathrm{EM}
= \sum_{\ell,\alpha}\hbar\omega_\ell\left(\hat{a}_{\ell,\alpha}^\dagger \hat{a}_{\ell,\alpha}+\frac{1}{2}\right).
\end{equation}
where we have used $[\hat{a}_{\ell,\alpha},\hat{a}_{\ell,\alpha}^\dagger] = 1$.
Introducing the canonical position and momentum as
\begin{equation}
\hat{q}_{\ell,\alpha}=\sqrt{\frac{\hbar}{2\omega_\ell}}\bigl(\hat{a}_{\ell,\alpha}^\dagger+\hat{a}_{\ell,\alpha}\bigr),
\qquad
\hat{p}_{\ell,\alpha}=\mathrm{i}\sqrt{\frac{\hbar\omega_\ell}{2}}\bigl(\hat{a}_{\ell,\alpha}^\dagger-\hat{a}_{\ell,\alpha}\bigr),
\end{equation}
we can recover the simple harmonic oscillator form:
\begin{equation}
\hat{H}_\mathrm{EM}
= \frac{1}{2}\sum_{\ell,\alpha}\Big(\hat{p}_{\ell,\alpha}^2+\omega_\ell^2\,\hat{q}_{\ell,\alpha}^2\Big).
\label{eq:HF-sho}
\end{equation}


Next, we simplify the molecular Hamiltonian, which still contains the vector potential field in the kinetic energy. 
Due to the separation between molecular and cavity length scales, we can approximate the vector potential as spatially uniform.
\begin{equation}
\hat{\bm A}(\bm r) \approx  \hat{\bm A}(0) =\sum_{\ell,\alpha}\frac{u_\ell(0)}{\sqrt{\epsilon_0}} \hat{q}_{\ell,\alpha}\hat{\bm e}_\alpha %\hat{\bm e}_\alpha \,.
\end{equation}
Here, we identify the so-called light-matter coupling strength as $\lambda_\ell = u_\ell(0)/\sqrt{\epsilon_0}$. %\frac{\epsilon_0}{\sqrt{ V}}$.
We now perform two unitary transformations to transfer the coupling from matter to photonic degrees of freedom.
The first eliminates the vector potential from the kinetic energy of matter by shifting it into the light momenta.
Let the total molecular dipole moment be $\hat{\bm \mu} = \sum_i Z_i e \hat{\bm R}_i - \sum_j e \hat{\bm r}_j$. We define
\begin{equation}
\hat{U}_L = \exp \left[\frac{\mathrm{i}}{\hbar}\hat{\bm \mu} \cdot \hat{\bm A}(0) \right] = \exp \left[\frac{\mathrm{i}}{\hbar} \sum_{\ell,\alpha} \lambda_\ell \,\hat{\mu}_\alpha \, \hat{q}_{\ell,\alpha} \right] 
\end{equation}
where $\hat{\mu}_\alpha = \hat{\bm\mu} \cdot \hat{\bm e}_\alpha$. The action of this transformation on the relevant operators follows from the Baker-Campbell-Hausdorff formula.
Using the commutation relations $[\hat{q}_{\ell,\alpha},\hat{p}_{\ell',\alpha'}]=i\hbar \delta_{\ell\ell'} \delta_{\alpha\alpha'}$ and $[\hat{\mu}_\alpha,\hat{q}_{\ell,\alpha'}]=0$, we obtain:
\begin{gather}
\hat{U}_{L}^{\dagger} \hat{p}_{\ell,\alpha} \hat{U}_{L} = \hat{p}_{\ell,\alpha} + \lambda_\ell \hat{\mu}_\alpha \,, \quad \hat{U}_{L}^{\dagger} \hat{q}_{\ell,\alpha} \hat{U}_{L} = \hat{q}_{\ell,\alpha} \\
\hat{U}_{L}^{\dagger} \hat{\bm P}_i \hat{U}_{L} = \hat{\bm P}_i + Z_i e\, \hat{\bm A}(0) \,, \quad \hat{U}_{L}^{\dagger} \hat{\bm p}_j\hat{U}_{L} = \hat{\bm p}_j - e \hat{\bm A}(0) \\
\hat{U}_{L}^{\dagger} \hat{\bm A}(0) \hat{U}_{L} = \hat{\bm A}(0) \,.
\end{gather}
Consequently, the canonical momenta transform as:
\begin{align}
\hat{U}_{L}^{\dagger}\, \big(\hat{\bm P}_i - Z_i e \hat{\bm A}(0)\big)\, \hat{U}_{L} &= \hat{\bm P}_i \\
\hat{U}_{L}^{\dagger}\, \big(\hat{\bm p}_j + e \hat{\bm A}(0)\big)\, \hat{U}_{L} &= \hat{\bm p}_j \,.
\end{align}
Applying $\hat{U}_{L}$ to the full Hamiltonian eliminates the vector potential from the matter kinetic energy and shifts the light momenta in $\hat{H}_\mathrm{EM}$:
\begin{align}
\hat{H}' \equiv \hat{U}_{L}^{\dagger} \hat{H} \hat{U}_{L}
&= \underbrace{\left[\sum_i \frac{\hat{\bm P}_i^{\,2}}{2M_i} + \sum_j \frac{\hat{\bm p}_j^{\,2}}{2m_j} + \hat{V}_{nn}+\hat{V}_{ee}+\hat{V}_{ne}\right]}_{\displaystyle \hat{H}_\mathrm{M}\ \text{(no vector potential)}} \notag\\
&\quad + \frac{1}{2} \sum_{\ell,\alpha}\bigg[\Big(\hat{p}_{\ell,\alpha} + \lambda_\ell \hat{\mu}_\alpha\Big)^{\!2} + \omega_\ell^2 \hat{q}_{\ell,\alpha}^2\bigg] .
\label{eq:HL-dE}
\end{align}

The second transformation moves the coupling from the photonic momenta to their positions via a $90^\circ$ phase rotation. A unitary transformation that achieves this is equivalent to time propagation of each cavity mode from $t=0$ to  $t=\pi/(2\omega_\ell)$:
\begin{equation}
\hat{U}_{\pi/2} \equiv
\exp\!\left[\frac{\mathrm i\,\pi}{4\hbar}\sum_{\ell,\alpha}\frac{\hat{p}_{\ell,\alpha}^{2}
+\omega_\ell^{2}\hat{q}_{\ell,\alpha}^{2}}{\omega_\ell}\right],
\end{equation}
which yields
\begin{equation}
\hat{U}_{\pi/2}^{\dagger}\,\hat{q}_{\ell,\alpha}\,\hat{U}_{\pi/2} = -\,\frac{\hat{p}_{\ell,\alpha}}{\omega_\ell},
\qquad
\hat{U}_{\pi/2}^{\dagger}\,\hat{p}_{\ell,\alpha}\,\hat{U}_{\pi/2} = \omega_\ell \hat{q}_{\ell,\alpha}.
\end{equation}
Acting with $\hat{U}_{\pi/2}$ on \cref{eq:HL-dE} and relabeling the rotated variables as $(\hat{q}_{\ell,\alpha},\hat{p}_{\ell,\alpha}) \to (\hat{p}_{\ell,\alpha},\hat{q}_{\ell,\alpha})$ gives
\begin{equation}
\hat{H}_{\mathrm{PF}}
= \hat{U}_{\pi/2}^{\dagger}\, \hat{H}'\, \hat{U}_{\pi/2}
= \hat{H}_\mathrm{M} + \frac{1}{2} \sum_{\ell,\alpha}\left[\hat{p}_{\ell,\alpha}^2
+ \omega_\ell^2\!\left(\hat{q}_{\ell,\alpha} + \frac{\lambda_\ell\hat{\mu}_\alpha}{\omega_\ell}\right)^{\!2} \right],
\label{eq:PF-final}
\end{equation}
which is the multi-mode Pauli--Fierz Hamiltonian in coordinate-displacement form where the molecular dipole moment displaces the light coordinate in some proportion.
In summary, the fully quantum light-matter Hamiltonian for the electron-nuclei-photonic system can be written as %We rewrite this compactly as:
\begin{gather}
\hat{H}_\mathrm{PF} = \hat{H}_\mathrm{M}+\hat{H}_\mathrm{EM} \,, \label{eq:H-final}
\\
\hat{H}_\mathrm{M} = \sum_i \frac{|\hat{\bm P}_i|^{2}}{2M_i} + \sum_j \frac{|\hat{\bm p}_j|^{2}}{2m_j} + \hat{V}_{nn}+\hat{V}_{ee}+\hat{V}_{ne} \,, \label{eq:HM-final}
\\
\hat{H}_\mathrm{EM} = \frac{1}{2} \sum_{\ell,\alpha}\left[
\hat{p}_{\ell,\alpha}^2 + \omega_\ell^2\!\left(\hat{q}_{\ell,\alpha} + \frac{\lambda_\ell \hat{\mu}_\alpha }{\omega_\ell}\right)^{2} \right] \,. \label{eq:HF-final}
\end{gather}

At this stage, we exploit timescale separations in the system.
Nuclei move more slowly than electrons, and for vibrational polaritons the EM resonance frequency typically matches nuclear vibrational frequencies or lies below them, placing both on similar timescales.
We therefore employ the cavity Born-Oppenheimer (BO) approximation, solving the electronic structure problem separately from the nuclear and photonic degrees of freedom:
\begin{gather}
\hat{H}_\mathrm{e} | \psi_g \rangle = V_\mathrm{BO}(\{ \bm{R}_i,q_{\ell,\alpha} \}) | \psi_g \rangle \,, \\
\hat{H}_\mathm{e} = \sum_j \frac{|\hat{\bm p}_j|^{2}}{2m_j}+\hat{V}_{nn}+\hat{V}_{ee}+\hat{V}_{ne} \,.
\end{gather}
Since the electromagnetic modes are off-resonant from electronic transitions, we assume they are decoupled from the electronic structure problem, so that $V_\mathrm{BO}(\{ \bm{R}_i,q_{\ell,\alpha} \}) \approx V_\mathrm{BO}(\{ \bm{R}_i\})$. %, avoiding the treatment of polaritonic potential energy surfaces.
Within this cavity BO approximation, treating the slower degrees of freedom classically, the potential energy of the classical light-matter system becomes:
\begin{gather}
V(\{ \bm{R}_i,q_{\ell,\alpha} \}) = V_\mathrm{BO}(\{ \bm{R}_i \})  + V_\mathrm{EM}(\{ \bm{R}_i,q_{\ell,\alpha} \})
\\
V_\mathrm{EM}(\{ \bm{R}_i,q_{\ell,\alpha} \}) =  \sum_{\ell,\alpha} \left[\frac{1}{2} \omega_\ell^2 q_{\ell,\alpha}^2+ \lambda_\ell \omega_\ell \langle \psi_g |\hat{\mu}_\alpha |\psi_g \rangle q_{\ell,\alpha} + \frac{\lambda_\ell^2}{2 } \langle \psi_g | \hat{\mu}_\alpha^2 |\psi_g \rangle \right] \,.
\end{gather}
The term $\langle \psi_g | \hat{\mu}_\alpha^2 |\psi_g \rangle$ represents quantum fluctuations of the dipole moment operator and, like $V_\mathrm{BO}$, also depends on $\{ \bm{R}_i,q_{\ell,\alpha} \}$.
For the classical treatment, we employ a mean-field approximation: $\langle \psi_g | \hat{\mu}_\alpha^2 |\psi_g \rangle \approx \left(\langle \psi_g | \hat{\mu}_\alpha |\psi_g \rangle \right)^2$, and defining the classical dipole moment as $\mu_\alpha(\{\bm{R}_i\}) = \langle \psi_g | \hat{\mu}_\alpha |\psi_g \rangle$, i.e., independent of photonic degrees of freedom, we can obtain %the following classical light-matter Hamiltonian:
\begin{gather}
H  = H_\mathrm{M}+H_\mathrm{EM} \label{eq:H-classical}
\\
H_\mathrm{M} = \sum_i \frac{|\bm{P}_i|^{2}}{2M_i} + V_\mathrm{BO}(\{ \bm{R}_i \}) \label{eq:HM-classical}
\\
H_\mathrm{EM} = \sum_{\ell,\alpha} \left[ \frac{p_{\ell,\alpha}^2}{2}+ \frac{1}{2} \omega_\ell^2 \left(q_{\ell,\alpha} +\frac{\lambda_\ell \mu_\alpha}{\omega_\ell} \right)^2 \right] \,,\label{eq:HF-classical}
\end{gather}
where $\mu_\alpha \equiv \mu_\alpha(\{ \bm{R}_i \})$.
\Cref{eq:H-classical,eq:HF-classical} provide a practical framework for investigating cavity effects through classical molecular dynamics (MD).
One can adopt models for the potential energy surface $V_\mathrm{BO}(\{ \bm{R}_i\})$ and dipole moment $\mu_\alpha(\{\bm{R}_i\})$ of increasing sophistication: non-polarizable empirical force fields, polarizable many-body force fields, machine learning models, or first-principles calculations.
In this work, we employ a simple non-polarizable classical force field.
A final simplification retains only the lowest-frequency EM mode, i.e., $\omega_{\ell=1} = \frac{\pi c }{L} = \omega_\mathrm{c}$.
This approximation is well justified by the energy-scale separation between relevant molecular vibrational processes and higher-order cavity modes, leading to the single-mode equations employed in this study.



\subsection{Modeling the Bath Environments}
\label{app:bath}

Hamilton's equations applied to \cref{eq:H-classical,eq:HF-classical} generate the deterministic equations of motion.
In a nonequilibrium setting, however, the light--matter system must be treated as open, exchanging energy with its surroundings at realistic rates.
Because the molecular subsystem and the cavity mode couple to physically distinct environments, each is assigned its own bath, as described below.


\subsubsection{Cavity Bath} 

Even in high-quality optical cavities, we still lose photons due to radiative losses.
We model this dissipation through coupling with a thermal electromagnetic environment outside the cavity. 
As is customary in modeling open systems, we model such an environment by coupling the cavity mode coordinate $q_\mathrm{c}$ (frequency $\omega_\mathrm{c}$, with  unit  mass) linearly to a continuum of bath oscillators.  % [Muhammad: Done - Added Caldeira-Leggett ref]
In the Caldeira--Leggett form~\cite{caldeira1983path}, the bath Hamiltonian can be written as
\begin{equation}
H_\mathrm{b} = \sum_{b} \left[\frac{p_b^2}{2} + \frac{1}{2} \omega_b^2 \left(x_b - \frac{c_b}{ \omega_b^2} q_\alpha\right)^2 \right]
\end{equation}
where $(x_b, p_b)$ are the bath oscillator coordinates with frequencies $\omega_b$ and coupling constants $c_b$. Integrating out the bath degrees of freedom yields the generalized Langevin equation
\begin{equation}
\ddot{q}_\alpha(t) + \omega_\alpha^2 q_\alpha(t) + \int_0^t \!\! dt' \, G(t-t') \, \dot{q}_\alpha(t') = \eta_\alpha(t) \,.
\end{equation}
Defining the spectral density as
\begin{equation}
J(\omega) = \frac{\pi}{2} \sum_{b} \frac{c_b^2}{\omega_b} \delta(\omega - \omega_b) \,,
\end{equation}
the memory kernel and noise correlation are related as follows:
\begin{gather}
G(t) = \frac{2}{\pi} \int_0^\infty \!\! d\omega \, \frac{J(\omega)}{\omega} \cos(\omega t) \,, \quad \langle \eta_\mathrm{c}(t) \eta_\mathrm{c}(t') \rangle = k_\mathrm{B} T \, G(t-t') \,.
\end{gather}
In the Ohmic, Markovian limit where $J(\omega) \approx \gamma_\mathrm{c} \omega$ at low frequencies, the memory kernel becomes $G(t) \approx 2\gamma_\mathrm{c} \delta(t)$, yielding standard Langevin dynamics,
\begin{gather}
\dot{q}_{\alpha} = p_\alpha \label{eq:langevin_cavity_1} \,, \\
\dot{p}_\alpha = -\omega_\mathrm{c}^2 q_\alpha - \lambda \omega_\mathrm{c} \mu_\alpha -\gamma_\mathrm{c} p_\alpha +\sqrt{2 \gamma_\mathrm{c} k_\mathrm{B} T } \eta_\alpha(t)  \,, \label{eq:langevin_cavity_2}
\end{gather}
where $\eta_\alpha(t)$ is white noise (the $\lambda=0$ case recovers the uncoupled cavity).
The lifetime of the cavity $\tau_\mathrm{c}$ is related to the damping rate through $\gamma_\mathrm{c} = 1/(2\tau_\mathrm{c})$.
When we turn on the coupling, $q_\mathrm{c} \to q_\mathrm{c} + \lambda \mu_\alpha / \omega_\mathrm{c}$, the damping and noise statistics remain unchanged while the deterministic force acquires the coupling term.

\subsubsection{Molecular bath.}

Experimentally, the thermal environment of a molecular system mainly mediates heat transfer between macroscopic boundaries and plays a limited role in structural relaxation. 
However, the simplest models for a stochastic bath, e.g., Langevin dynamics, act locally on nuclear degrees of freedom, which can severely distort intrinsic dynamics if heat transfer proceeds fast. %because the bath acts locally, perturbing each nucleus with independent white noise. 
The Bussi--Parrinello (BP) thermostat\cite{bussi2007canonical,bussi2008stochastic} solves this problem by thermalizing the system like a Langevin bath while minimizing local perturbations to the dynamics. 
The BP bath is derived by analyzing Langevin dynamics, where each particle experiences independent friction and noise through
\begin{gather}
\dot{\bm R}_i = \frac{\bm{P}_i}{M_i}, \\
\dot{\bm{P}}_i = -\nabla_{\bm{R}_i} V - \gamma_\mathrm{b} \bm{P}_i + \sqrt{2M_i\gamma_\mathrm{b} k_\mathrm{B}T}\, \bm{\eta}_i(t) \,,
\end{gather}
where $\bm{\eta}_i(t)$ are independent white noise processes with 
$\langle {\bm \eta}_i(t) {\bm \eta}_j(t') \rangle = {\bm \delta}_{ij} \delta(t-t')$.  
The Bussi--Parrinello scheme exploits the key observation that, if we calculate the time derivative of the total energy $H(t)$ using the Itô chain rule and sum over all noise terms, we obtain
\begin{equation}
\dot{H}(t) = -\frac{K(t) - \bar{K}}{\tau_\mathrm{b}} + 2\sqrt{\frac{K(t)\bar{K}}{3N\tau_\mathrm{b}}} \, \eta(t) \,,
\label{eq:energyrate_langevin}
\end{equation}
where $\bar{K} = \tfrac{3}{2}Nk_\mathrm{B}T$, $\tau_\mathrm{b} = (2\gamma_\mathrm{b})^{-1}$, and $\eta(t)$ is a single global white noise emerging from the sum of $3N$ independent noise terms. 
Since only the total kinetic energy $K$ appears in the energy evolution, we can design an alternative stochastic force $\bm{G}_i$ that depends globally on $K$ and reproduces the same rate of energy exchange while minimizing perturbations to individual particle trajectories. 
The minimization step follows from quantifying the disturbance on the kinetic energy as\cite{bussi2008stochastic}
\begin{equation}
D = \sum_{i=1}^{N} \frac{|\bm G_i|^2}{2M_i}
\end{equation}
and minimizing it subject to the constraint that $\dot{H}(t)$ matches the Langevin result. 
Geometrically, the constraint forces $\bm G_i$ to be parallel to each momentum vector, i.e., 
\begin{equation}
\bm{G}_i = \bm{P}_i \left[-\gamma_K + \sqrt{2D_K}\,\eta(t)\right],
\end{equation}
where the scalar functions $\gamma_K$ and $D_K$ depend on the total kinetic energy $K$ and are determined by matching the energy evolution. 
The change in total energy under this force is
\begin{equation}
\dot{H} = -2\gamma_K K + 2D_K K + 2\sqrt{2D_K}\,K\,\eta(t).
\end{equation}
Enforcing the equality with \cref{eq:energyrate_langevin} gives
\begin{gather}
\gamma_K = \frac{1}{2\tau_\mathrm{b}}\left(1-\left(1-\frac{1}{3N}\right)\frac{\bar{K}}{K}\right) \,, \\
D_K = \frac{\bar{K}}{6N\,\tau_\mathrm{b}\,K} \,.
\end{gather} 
The equations of motion thus become
\begin{gather}
\dot{\bm{R}}_i = \frac{\bm{P}_i}{M_i},  \\
\dot{\bm{P}}_i = -\nabla_{\bm{R}_i} V -\gamma_K \bm P_i + \sqrt{2 D_K} \bm P_i \eta(t). 
\end{gather}
This thermostat samples the canonical ensemble exactly and preserves short-time velocity autocorrelation functions, which are critical for maintaining correct microscopic dynamics in supercooled liquids.\cite{bussi2008stochastic} 
For a single degree of freedom ($N=1$), it reduces to standard Langevin dynamics, whereas in many-particle systems, the global coupling reduces dynamical artifacts coming from the bath coupling.


\section{Observables in an Aging System} \label{app:observables}


\subsection{Averaged Observables} \label{app:avgobserv}

During an aging experiment, the main prbe of dynamics is the ensemble average of an observable 
$A(\Gamma)$ at time $t$, e.g., the average potential energy or fictive temperature. Here, $\Gamma=(\bm R, \bm P, \bm q, \bm p)$ denotes a phase-space point of our light-matter system, with $\bm R = (\bm R_1,\ldots,\bm R_N)$, $\bm P = (\bm P_1,\ldots,\bm P_N)$, $\bm q = (q_1,q_2)$, and $\bm p = (p_1,p_2)$. Let the phase-space distribution at time $t$ be $\rho^{(\lambda)}(\Gamma,t)$; the ensemble average of $A(\Gamma)$ at time $t$ is
\begin{equation}
\langle A(t) \rangle_\lambda = \int \mathrm{d}\Gamma\, A(\Gamma)\,
\rho^{(\lambda)}(\Gamma, t) \,.
\end{equation}
Setting the time at which we activate the coupling to be $t_0$, we can write the Fokker--Planck equation\cite{risken1996fokker} for our non-thermal aging protocol as
\begin{equation}
\frac{\partial \rho^{(\lambda)}(\Gamma,t)}{\partial t} 
= \mathcal{L}_\lambda\,\rho^{(\lambda)}(\Gamma,t) \,,
\qquad
\mathcal{L}_\lambda= \mathcal{L}_0 + \mathcal{L}_\mathrm{int}(\lambda) + \mathcal{L}_\mathrm{b},
\end{equation}
with initial condition $\rho^{(\lambda)}(\Gamma,t_0) = \rho_\mathrm{eq}^{(\lambda=0)}(\Gamma)$, and the equilibrium distribution
\begin{equation}
\rho_\mathrm{eq}^{(\lambda)}(\Gamma) = \frac{1}{Z(T, \lambda)}e^{-\beta H(\Gamma;\lambda)} \,, 
\qquad 
Z(T, \lambda) = \int \mathrm{d} \Gamma \, e^{-\beta H(\Gamma;\lambda)} \,,
\end{equation}
where the total Hamiltonian $H(\Gamma; \lambda)$ is equivalent to \cref{eq:H-classical}. Note that the Fokker-Planck equation already models the step change in the coupling constant by virtue of starting with an initial condition at zero coupling.  

The operator $\mathcal{L}_\lambda$ consists of the bare part $\mathcal{L}_0$, interaction part $\mathcal{L}_\mathrm{int}(\lambda)$, and bath contributions $\mathcal{L}_\mathrm{b}$. The bare part $\mathcal{L}_0$ is the Liouvillian operator generated by the original Hamiltonian dynamics at zero coupling. Its action on an arbitrary distribution $\rho$ is
\begin{equation}
\mathcal L_0 \rho
= - \sum_{i=1}^N \frac{\bm P_i}{M_i} \cdot\nabla_{\bm R_i} \rho
  + \sum_{i=1}^N \nabla_{\bm R_i} V(\bm R) \cdot \nabla_{\bm P_i} \rho
  -\sum_{\alpha=1,2} p_\alpha \frac{\partial \rho}{\partial q_\alpha}
  + \sum_{\alpha=1,2} \omega_\mathrm{c}^2 q_\alpha \frac{\partial \rho}{\partial p_\alpha} \,.
\end{equation}
The interaction part $\mathcal{L}_\mathrm{int}(\lambda)$ arises from the light--matter coupling term:
\begin{equation}
\mathcal L_{\mathrm{int}}(\lambda)\,\rho
= \sum_{i=1}^N  
\left[
   \sum_{\alpha=1,2} \lambda\,z_i e
   \bigl(\omega_\mathrm{c} q_\alpha + \lambda\,\mu_\alpha(\bm R)\bigr)
   \hat{\bm e}_\alpha
\right] \cdot \nabla_{\bm P_i} \rho
 - \sum_{\alpha=1,2} \lambda\,\omega_\mathrm{c}\,\mu_\alpha(\bm R)\,
   \frac{\partial \rho}{\partial p_\alpha}.
\end{equation}
Both noise and dissipation terms from the molecular and cavity baths are lumped into $\mathcal{L}_\mathrm{b}$, which we write as
\begin{equation}
\mathcal L_\mathrm{b} = \mathcal L_\mathrm{L} + \mathcal L_{\mathrm{BP}} \,,
\end{equation}
where the cavity Langevin (Kramers) operator\cite{kramers1940brownian,risken1996fokker} acting on the cavity momenta is
\begin{equation}
\mathcal L_\mathrm{L}\,\rho
= \sum_{\alpha=1,2} \left[
\frac{\partial}{\partial p_\alpha}\bigl(\gamma_\mathrm{c} p_\alpha\,\rho\bigr)
+ \gamma_\mathrm{c} k_\mathrm{B} T\,
  \frac{\partial^2 \rho}{\partial p_\alpha^2}
\right] \,,
\end{equation}
and the Bussi--Parrinello operator acts only on the particle momenta with a different diffusivity tensor,
\begin{equation}
\mathcal L_{\mathrm{BP}}\,\rho
= \sum_{i=1}^N \nabla_{\bm P_i}\!\cdot\!\bigl(\gamma_K\,\bm P_i\,\rho\bigr)
  + \sum_{i=1}^N \sum_{j=1}^N \bigl(\nabla_{\bm P_i}\otimes\nabla_{\bm P_j}\bigr)
    : \left[ D_K \left(\bm P_i\otimes\bm P_j\right) \rho \right]\,,
\label{eq:LBP_compact}
\end{equation}
with $\otimes$ denoting the outer product.

Delegating the time dependence into the phase-space distribution is known as the Schr\"odinger picture of the dynamics.\cite{tuckerman2010statistical} An equivalent way to average the observable is via the Heisenberg picture, where given an initial condition $\Gamma$, we directly evolve the observable averaged over all possible realizations of the noise, i.e., 
\begin{equation}
A_{t}(\Gamma; \lambda) = \mathbb{E}\left[ A(\Gamma(t; \lambda)) \, | \, \Gamma(0)=\Gamma \right] \,.
\end{equation}
We then write the equivalent ensemble average as
\begin{equation}
\langle A(t) \rangle_\lambda = \int \mathrm{d}\Gamma \, A_{t}(\Gamma; \lambda) \,
\rho^{(\lambda=0)}_\mathrm{eq}(\Gamma) \,.
\end{equation}
The conditionally averaged observable $A_{t}(\Gamma; \lambda)$ obeys the adjoint version of the Fokker-Planck equation, also known as the backward Kolmogorov equation:\cite{risken1996fokker,gardiner2009stochastic}
\begin{equation}
\frac{\partial A_{t}(\Gamma; \lambda)}{\partial t} = \mathcal{L}^\dagger_\lambda A_t(\Gamma; \lambda) \,.
\end{equation}
with the initial condition $A_{t=0}(\Gamma; \lambda) = A(\Gamma)$.

When we run molecular simulations, we implicitly assume a Heisenberg picture by taking a sample mean over $N_\mathrm{T}$ independent trajectories, initiated from a sample of initial configurations at equilibrium. Denoting $A^{(n)}(t;\lambda) = A(\Gamma^{(n)}(t;\lambda))$ as the observable evaluated at the phase-space point $\Gamma^{(n)}(t;\lambda)$ at time $t$ from the $n$-th trajectory, we obtain
\begin{equation}
\langle A(t) \rangle_\lambda \approx \frac{1}{N_\mathrm{T}} \sum_{n=1}^{N_\mathrm{T}} A^{(n)}(t;\lambda) \,.
\end{equation}
%Since the EOMs generate the trajectory, we fix it to be one realization of the noise per initial configuration. 
The averaging over independent trajectories is really an average over all possible initial conditions $\Gamma^{(n)}_0$, which are \textit{a priori} sampled from another molecular simulation.

When calculating \textit{equilibrium averages}, we note that the system reaches its new equilibrium, and we can write this ensemble average as 
\begin{equation}
\langle A\rangle_\lambda  = \int \mathrm{d} \Gamma\, A(\Gamma)\, \rho^{(\lambda)}_\mathrm{eq}(\Gamma) \,.
\end{equation}
An important assumption we make regarding systems at equilibrium is the \textit{ergodic hypothesis},\cite{tuckerman2010statistical,hansen2013theory} i.e., the equivalence of the ensemble average and the long-time average over a single noisy trajectory:
\begin{equation}
\langle A\rangle_\lambda  = \bar{A}(\Gamma_0; \lambda) 
= \lim_{T \to \infty} \frac{1}{T} \int_0^T \mathrm{d}t\, A(\Gamma(t;\lambda)) \,.
\end{equation}
The practical utility of $\langle A\rangle_\lambda$ is to supplement statistical averaging of our time averaging. First, for the $n$-th trajectory, we define its time average by subsampling at discrete times $t_p = p \Delta t$, with $p = 0,1,\ldots, M-1$, 
\begin{equation}
\bar{A}(\Gamma_0^{(n)}; \lambda) \approx \frac{1}{M} \sum_{p=0}^{M-1} A\bigl(\Gamma^{(n)}(t_p;\lambda)\bigr) \,,
\end{equation}
and thus we approximate the equilibrium average as
\begin{equation}
\langle A \rangle_\lambda \approx \frac{1}{N_\mathrm{T}} \sum_{n=1}^{N_\mathrm{T}}  \bar{A}(\Gamma_0^{(n)}; \lambda) 
\approx \frac{1}{M N_\mathrm{T}} \sum_{n=1}^{N_\mathrm{T}} \sum_{p=0}^{M-1} A\bigl(\Gamma^{(n)}(t_p;\lambda)\bigr) \,. \label{eq:Aeq_samplemean}
\end{equation}

In an aging study, we typically denote the time elapsed after the beginning of the protocol as the waiting time $t_\mathrm{w} = t-t_0$. For example, suppose we choose density-mode fluctuations at a wavevector $\bm k$, defined here as 
\begin{equation}
\rho_{\bm k}^{(n)}(t; \lambda) = \sum_{i=1}^N e^{\mathrm{i}\bm k \cdot \bm R_i^{(n)}(t; \lambda)} \,,
\end{equation}
where $\bm R_i^{(n)}(t;\lambda)$ is the $i$-th atom position at time $t$ for non-zero $\lambda$. The so-called static structure factor is defined as the second moment of the density fluctuations
\begin{equation}
S^{(\lambda)}_{\bm k}(t_\mathrm{w}) = \left\langle \rho_{\bm k}(t_\mathrm{w})  \rho_{\bm k}^*(t_\mathrm{w})  \right\rangle_\lambda \approx  \frac{1}{N_\mathrm{T}} \sum_{n=1}^{N_\mathrm{T}} \rho_{\bm k}^{(n)}(t_\mathrm{w};\lambda)  \rho_{\bm k}^{*(n)}(t_\mathrm{w};\lambda) \,,
\end{equation}
The term `static' refers to the fact that, at long waiting times where the system converges to equilibrium, the average becomes independent of time, and we can make use of \cref{eq:Aeq_samplemean} to obtain %of
\begin{equation}
S^{(\lambda)}_{\bm k} = \left\langle \rho_{\bm k}  \rho_{\bm k}^* \right\rangle_\lambda \approx  \frac{1}{N_\mathrm{T} M} \sum_{n=1}^{N_\mathrm{T}} \sum_{p=0}^{M-1}\rho_{\bm k}^{(n)}(t_p;\lambda)\rho_{\bm k}^{*(n)}(t_p;\lambda) 
\end{equation}

\subsection{Time-Correlation Functions} \phantomsection \label{app:timecorrel}

\subsubsection{General Case}
Central to the dynamic response of an aging system are its time-correlation functions, which for a general observable $A(t)$ take the form
\begin{equation}
C_{AA}^{(\lambda)}(t; t_\mathrm{w}) = \langle A(t+t_\mathrm{w})\,A(t_\mathrm{w}) \rangle_\lambda \,.
\label{eq:NEQ_CAA_def}
\end{equation}
In terms of the conditionally averaged observable at time $t+t_\mathrm{w}$, $A_{t+t_\mathrm{w}}(\Gamma)$,  and the phase-space distribution $\rho^{(\lambda)}(\Gamma,t_\mathrm{w})$ at waiting time $t_\mathrm{w}$, the ensemble average in the time-correlation function can be written
\begin{equation}
C_{AA}^{(\lambda)}(t;t_\mathrm{w})
= \int \mathrm d\Gamma\,
A_{t+t_\mathrm{w}}(\Gamma; \lambda)\,A(\Gamma)\,
\rho^{(\lambda)}(\Gamma,t_\mathrm{w}) \,.
\label{eq:CAA_ensavg}
\end{equation}
We use our cavMD simulations to approximate \cref{eq:CAA_ensavg} via sampling over many trajectories. In particular, we can perform ensemble averaging by generating $N_\mathrm{T}$ trajectories, simulating from equilibrium configurations at zero coupling at $t = t_0$ directly up to $t+t_\mathrm{w}$. We then use the two points in the $n$-th trajectory to calculate the time-correlation function,
% \footnote{This is slightly different from following the definition of \cref{eq:CAA_ensavg}, which would imply the following two steps. First, for a given waiting time $t_\mathrm{w}$ we simulate starting from equilibrium configurations at zero coupling and time $t = t_0$ until time $t =t_\mathrm{w}$; the final configuration at time $t=t_\mathrm{w}$ samples the distribution $\rho^{(\lambda)}(\Gamma,t_\mathrm{w})$. 
% Second, we perform conditional averaging by taking each such configuration from $t=t_\mathrm{w}$ and running multiple trajectories from it until $t = t_\mathrm{w}$.
% Because the noises are statistically independent, both approaches are equivalent. 
% }
\begin{equation}
C_{AA}^{(\lambda)}(t;t_\mathrm{w}) \approx \frac{1}{N_\mathrm{T}} \sum_{n=1}^{N_\mathrm{T}}
A^{(n)}(t+t_\mathrm{w}; \lambda)\, A^{(n)}(t_\mathrm{w}; \lambda) \,.
\end{equation}
Using the density-mode fluctuations again as an example, we compute their time-correlation function $F_{\bm k}(t;t_\mathrm{w})$, the intermediate scattering function (ISF).\cite{hansen2013theory} Averaging over $N_{\bm k}$ wavevectors of equal magnitude $|\bm k|$ distributed on a uniformly gridded sphere yields
\begin{equation}
F_{\bm k}^{(\lambda)}(t;t_\mathrm{w}) \approx  \frac{1}{N_\mathrm{T} N_{\bm k}} \sum_{n=1}^{N_\mathrm{T}} \sum_{s=1}^{N_{\bm k}}
\rho_{\bm k_s}^{(n)}(t+t_\mathrm{w}; \lambda)\,\rho_{\bm k_s}^{*(n)}(t_\mathrm{w}; \lambda) \,.
\label{eq:Fk_NEQ_fixed}
\end{equation}
and the normalized correlator $\phi_{k}(t;t_\mathrm{w}) = F_{k}(t;t_\mathrm{w})/S_{k}(t_\mathrm{w})$. Note that the averaging over wavevectors is also known as $\bm k$-averaging. 
To track nonequilibrium dynamics, we maintain multiple waiting times $t_\mathrm{w}^{(r)}$ in a single simulation, spaced at regular intervals. At each timestep, we compute $F_{k}(t;t_\mathrm{w}^{(r)})$ at a single trajectory level. For reporting, we use $N_\mathrm{T}=500$ trajectories and $N_{\bm k} = 50$ wavevectors. %per $k$-shell. %When desired, we also present to compare across waiting times and state points.

\subsubsection{Equilibrium Case}
% %The right-hand side depends only on the time difference $t$ between the two observations and not on the absolute waiting time $t_\mathrm{w}$. We therefore define the equilibrium autocorrelation function
% \begin{equation}
% C_{AA,\mathrm{eq}}^{(\lambda)}(t)
% := C_{AA}^{(\lambda)}(t;t_\mathrm{w}) 
% = \int \mathrm d\Gamma\, \rho^{(\lambda)}_\mathrm{eq}(\Gamma)\,
% A(\Gamma)\,\bigl(e^{t\mathcal L^\dagger} A\bigr)(\Gamma) \,.
% \end{equation}

% It is often convenient to express the same quantity using 
In the limit of large waiting times, the system reaches equilibrium (for fixed $\lambda$) and time-correlation functions obey time-translational invariance. A short proof of this property follows from \cref{eq:CAA_ensavg}. In this regime, we assume that at that time, the system has relaxed to the stationary distribution so that
\begin{equation}
\rho^{(\lambda)}(\Gamma,t_\mathrm{w}) = \rho^{(\lambda)}_\mathrm{eq}(\Gamma),
\qquad
\mathcal L_\lambda\,\rho^{(\lambda)}_\mathrm{eq} = 0 \,,
\end{equation}
which turns \cref{eq:CAA_ensavg} into %, and substitute the equilibrium distribution:
\begin{equation}
C_{AA}^{(\lambda)}(t;t_\mathrm{w})
= \int \mathrm d\Gamma A_{t+t_\mathrm{w}}(\Gamma) A(\Gamma) \,
\rho^{(\lambda)}_\mathrm{eq}(\Gamma) \,.
\end{equation}
To eliminate the last remaining dependence in $t_\mathrm{w}$, we note that $A_{t+t_\mathrm{w}}(\Gamma)$ is obtained from solving the backward Kolmogorov equation with an initial condition at $t =t_\mathrm{w}$. Since the equation is invariant under time translation, $A_{t+t_\mathrm{w}}(\Gamma) \to A_{t}(\Gamma)$ provided that the initial condition is set at $t = 0$. Thus, we obtain
\begin{equation}
C_{AA}^{(\lambda)}(t;t_\mathrm{w})
=  \int \mathrm{d} \Gamma A_{t}(\Gamma; \lambda) A(\Gamma) \rho_\mathrm{eq}^{(\lambda)}(\Gamma)
=: C_{AA,\mathrm{eq}}^{(\lambda)}(t), \label{eq:equilibrium_autocorr_ens}
\end{equation}
which makes explicit that the equilibrium correlation function depends only on the time difference $t$, i.e., it is time-translationally invariant.

At equilibrium, the ergodic hypothesis also holds for time-correlation functions. In other words, the ensemble average in \cref{eq:equilibrium_autocorr_ens} can be replaced by a long-time average along a single stationary trajectory $\Gamma(t;\lambda)$ given a fixed initial condition $\Gamma_0$:
\begin{equation}
C_{AA,\mathrm{eq}}^{(\lambda)}(\tau)
= \bar{C}_{AA}^{(\lambda)}(\tau;\Gamma_0) := \lim_{T \to \infty} \frac{1}{T} \int_0^T \mathrm{d}t'\,
A\bigl(\Gamma(\tau+t^\prime; \lambda) \bigr) A\bigl(\Gamma(t^\prime; \lambda) \bigr)
\label{eq:equilibrium_autocorr}
\end{equation}
In practice, one would typically combine the time averaging with ensemble averaging of the initial condition, i.e., sample $\Gamma(0;\lambda) = \Gamma_0  \sim \rho_{\mathrm{eq}}^{(\lambda)}$, to increase statistical convergence further, and thus
\begin{equation}
\bar{C}_{AA}^{(\lambda)}(\tau) = \int \mathrm{d} \Gamma_0 \ \bar{C}_{AA}^{(\lambda)}(\tau;\Gamma_0) \ \rho^{(\lambda)}_\mathrm{eq}(\Gamma_0)
\end{equation}
In our simulations, observables are sampled at discrete times $t_p = p\,\Delta t$ along each trajectory, with $p = 0,1,\ldots,M-1$. The equilibrium autocorrelation at lag $\tau_k = k \Delta t$ is then estimated as
\begin{equation}
\bar{C}_{AA}^{(\lambda)}(\tau_k)
\approx  \frac{1}{N_\mathrm{T}(M-k)} \sum_{n=1}^{N_\mathrm{T}} \sum_{p=0}^{M-k-1}
A^{(n)}(t_p;\lambda)\, A^{(n)}(t_{p+k};\lambda),
\end{equation}
with $k = 0, 1, \ldots, M-1$, where the factor $(M-k)^{-1}$ accounts for the decreasing number of available time pairs at larger lags. This discrete estimator is the standard “sliding time-origin” implementation of \cref{eq:equilibrium_autocorr}. % and does not rely on fast Fourier transforms.
Applying the same procedure to the density modes yields the equilibrium density correlation function
\begin{equation}
\bar{F}_{\bm k}(\tau_k)
\approx \frac{1}{N_\mathrm{T}N_{\bm k}(M-k)} \sum_{n=1}^{N_\mathrm{T}} \sum_{p=0}^{M-k-1}
\sum_{s=1}^{N_{\bm k}} \rho_{\bm k_s}^{(n)}(t_p)\, \rho_{\bm k_s}^{(n)*}(t_{p+k}) ,
\end{equation}
where the $\bm k$-vectors in a given shell have equal magnitude $|\bm k|$ and are distributed on a Fibonacci sphere as described above. To quantify the structural relaxation timescale from the normalized ISF $\phi_{\bm k}^{(\lambda)}(t;t_\mathrm{w})$ data, we define it as the time when $F_{\bm k}$ crosses a threshold value $a$ (typically $a = 0.1$):
\begin{equation}
\phi_{\bm k}^{(\lambda)}(\tau_\mathrm{s};t_\mathrm{w}) = a.
\label{eq:threshold_criterion}
\end{equation}
We locate $\tau_\mathrm{s} \equiv \tau_\mathrm{s}(t_\mathrm{w},\lambda)$ by linear interpolation between adjacent data points that bracket the threshold. 
Note that these relaxation times, extracted from $F_{\bm k}(t;t_\mathrm{w})$ at multiple waiting times $t_\mathrm{w}$, provide the data needed to construct the material time $h_\lambda(t)$, allowing us to quantify the effective aging rate under cavity coupling; see \cref{app:mattime} for details. 


For the dipole-dipole autocorrelation function (DACF), we subtract the mean before computing correlations to remove statistical bias in our estimate, since we know \textit{a priori} that the ensemble average should be exactly zero. Defining the fluctuations $\delta \mu_\alpha(t) = \mu_\alpha(t) - \bar{\mu}_\alpha$ where $\bar{\mu}_\alpha = M^{-1}\sum_{p=0}^{M-1} \mu_\alpha(t_p)$, the centered correlation function becomes
\begin{equation}
\bar{C}_{\mu \mu}^{(\lambda)}(\tau_k) \approx \frac{1}{N_\mathrm{T}(M-k)} \sum_{n=1}^{N_\mathrm{T}} \sum_{p=0}^{M-k-1}  \sum_{\alpha =1,2} \delta \mu_\alpha(t_p;\lambda) \, \delta \mu_\alpha(t_{p+k};\lambda),
\end{equation}
where the sum runs over the two polarization directions $\alpha$. For the equilibrium DACF, one important application is to calculate the IR absorption spectrum $\alpha(\omega)$ from the DACF. The abtion is related to the dipole-dipole correlation via
\begin{equation}
n(\omega) \alpha(\omega) = \frac{\beta \omega^2}{4 \epsilon_0 V c} \int_{-\infty}^{+\infty} \!\!\!\! dt \, e^{-i \omega t} C_{\mu\mu}(t)  \,, \label{eq:irspectrum}
\end{equation}
from which we note that
\begin{equation}
\int_{-\infty}^{\infty} d\tau \, e^{i \omega \tau}\, \bar{C}_{\mu \mu}^{(\lambda)}(\tau)
\to 2 \int_{0}^{\infty} d\tau \, \cos(\omega \tau)\, \bar{C}_{\mu \mu}^{(\lambda)}(\tau)
\end{equation}
which follows from time-reversibility, $\bar{C}_{\mu \mu}^{(\lambda)}(\tau)=\bar{C}_{\mu \mu}^{(\lambda)}(-\tau)$ and the fact we only care about the real part of the absorption spectrum. Truncating \cref{eq:irspectrum} at the final simulation time, $\tau_{\max}$ then leads to
\begin{equation}
n(\omega) \alpha(\omega) \approx \frac{\beta \omega^{2}}{4 \epsilon_0 V c} \int_{0}^{\tau_{\max}} \!\!\!\!\!\mathrm{d}\tau\,\cos(\omega \tau)\,\bar{C}_{\mu \mu}^{(\lambda)}(\tau) .
\end{equation}
We sample the DACF uniformly with spacing $\Delta t$ and $N$ points so that
$\tau_{\max} = (N-1) \Delta t$ and $C_n = C_\mu(n\,\Delta t)$. The frequencies are gridded according to
\begin{equation}
\omega_k = \frac{\pi k}{(N-1)\,\Delta t} .
\end{equation}
Using the discrete cosine transform, we compute
\begin{equation}
X_k = C_0 + (-1)^k C_{N-1} + 2 \sum_{n=1}^{N-2} C_n \cos\!\left(\frac{\pi n k}{N-1}\right) .
\end{equation}
The trapezoidal rule then approximates the cosine integral as
\begin{equation}
\int_{0}^{\tau_{\max}} d\tau \,\cos(\omega_k \tau)\, C_\mu(\tau) \approx \frac{\Delta t}{2}\,X_k ,
\end{equation}
yielding the discrete estimator for the IR spectrum
\begin{equation}
n(\omega_k)\,\alpha(\omega_k) \approx \frac{\beta\,\omega_k^{2}}{4\,\epsilon_0\,V\,c}\,\frac{\Delta t}{2}\,X_k .
\end{equation}
Note that when plotting intensities, we often normalize the spectrum by the area under the curve.


\subsection{Material Time} \phantomsection \label{app:mattime}

The Methods section summarizes the regularized least-squares formulation for reconstructing the material time $h_\lambda(t)$. Here we provide the full derivation, including the explicit constraint equations, the smoothness-penalty matrix, and the gradient calculation leading to the normal equations.

\subsubsection{Grid Parameterization and Constraints}

We discretize time on a dense regular grid of $M_\mathrm{grid}$ points: $t_m = m \cdot \Delta t$ for $m = 0, 1, \ldots, M_\mathrm{grid}-1$, where $\Delta t = t_{\max}/(M_\mathrm{grid}-1)$.
The parameter vector $\boldsymbol{h}_\lambda = [h_{\lambda}^0, h_{\lambda}^1, \ldots, h_{\lambda}^{M_\mathrm{grid}-1}]^\mathsf{T}$ stores the material-time values at each gridpoint, so that $\tilde{h}_\lambda(t_m) = h_\lambda^m$.
In practice, $M_\mathrm{grid}$ should be large enough (e.g., $M_\mathrm{grid} = 500$) that the discrete grid approximates continuous $h_\lambda(t)$ well, enabling smooth reconstructions via linear interpolation between gridpoints.

Substituting the piecewise linear expansion into the MTTI constraints, each constraint becomes
\begin{equation}
\sum_{m = 0}^{M_\mathrm{grid}-1}  \left( \theta_m(t_\mathrm{w}^{(r)} + \tau_\mathrm{s}^{(r)})  - \theta_m(t_\mathrm{w}^{(r)}) \right) h_\lambda^m= 1, \quad r = 1, 2, \ldots, N_\mathrm{w},
\label{eq:SI_constraints_set}
\end{equation}
which defines the $N_\mathrm{w} \times M_\mathrm{grid}$ interpolation-weight matrix $\mathbf{A}$ introduced in the Methods.
The unconstrained least-squares solution minimizes $\|\mathbf{A}\boldsymbol{h}_\lambda - \boldsymbol{b}\|^2$; however, because the system is heavily overparameterized ($M_\mathrm{grid} \gg N_\mathrm{w}$), this solution will overfit the data.

\subsubsection{Smoothness Regularization}

To prevent overfitting, we penalize solutions with large curvature via a Tikhonov term.
We define a $(M_\mathrm{grid}-2) \times M_\mathrm{grid}$ second-order finite-difference matrix $\mathbf{D}$ that computes the centered second derivative of $h_\lambda(t)$ on the grid, with elements
\begin{equation}
D_{rm} = 
\begin{cases}
1, & m = r \\
-2, & m = r+1 \\
1, & m = r+2 \\
0, & \text{otherwise}
\end{cases}
\end{equation}
for $r = 0, 1, \ldots, M_\mathrm{grid}-3$ and $m = 0, 1, \ldots, M_\mathrm{grid}-1$. Explicitly, $\mathbf{D}$ has the tridiagonal banded form:
\begin{equation}
\mathbf{D} = \begin{bmatrix}
1 & -2 & 1 & 0 & 0 & \cdots & 0 \\
0 & 1 & -2 & 1 & 0 & \cdots & 0 \\
0 & 0 & 1 & -2 & 1 & \cdots & 0 \\
\vdots & \vdots & \vdots & \ddots & \ddots & \ddots & \vdots \\
0 & 0 & 0 & \cdots & 1 & -2 & 1
\end{bmatrix}.
\end{equation}

\subsubsection{Derivation of the Normal Equations}

The regularized objective introduced in the Methods is
\begin{equation}
J(\boldsymbol{h}_\lambda)
= \|\mathbf{A}\boldsymbol{h}_\lambda - \boldsymbol{b}\|^2
+ \alpha \|\mathbf{D}\boldsymbol{h}_\lambda\|^2,
\end{equation}
where $\alpha > 0$ controls the smoothness--data tradeoff.
Taking the gradient with respect to $\boldsymbol{h}_\lambda$ and setting it to zero:
\begin{equation}
\nabla_{\boldsymbol{h}_\lambda} J
= 2\mathbf{A}^\mathsf{T}(\mathbf{A}\boldsymbol{h}_\lambda^* - \boldsymbol{b})
+ 2\alpha\,\mathbf{D}^\mathsf{T}\mathbf{D}\boldsymbol{h}_\lambda^*
= \mathbf{0}.
\end{equation}
Dividing by 2 and rearranging yields the normal equations
\begin{equation}
\bigl(\mathbf{A}^\mathsf{T}\mathbf{A}
+ \alpha\,\mathbf{D}^\mathsf{T}\mathbf{D}\bigr)\,\boldsymbol{h}_\lambda^*
= \mathbf{A}^\mathsf{T}\boldsymbol{b}.
\end{equation}
The matrix $\mathbf{A}^\mathsf{T}\mathbf{A} + \alpha\,\mathbf{D}^\mathsf{T}\mathbf{D}$ is symmetric positive-definite for $\alpha > 0$, and inherits a banded structure from $\mathbf{D}^\mathsf{T}\mathbf{D}$, so the system can be efficiently solved using sparse Cholesky decomposition or conjugate gradient methods.

\section{Simulation Details} \label{app:simulationdetails}

\subsection{Numerical Integration Schemes}

A summary of the numerical methods is provided in the Methods section of the main text. Here we present the full integration schemes.
All cavMD simulations employ split-operator integration schemes\cite{tuckerman2010statistical} that separate deterministic equations of motion from stochastic thermostatting.
The cavity mode follows Langevin dynamics integrated via a modified velocity-Verlet scheme. 
For the cavity mode with canonical coordinates $(q_\mathrm{\alpha}, p_\mathrm{\alpha})$ and unit mass, the first half-step advances the momentum by
\begin{equation}
p_\alpha\left(t + \Delta t/2\right) = p_\alpha(t) + \frac{1}{2} f_\alpha(t) \Delta t \,,
\end{equation}
where $f_\alpha(t) = -\omega_\alpha^2\!\left(q_\alpha + \lambda \mu_\alpha / \omega_\mathrm{c}\right)$ is the effective force that combines the harmonic restoring potential and the coupling term.
The position then updates according to the updated momenta.
\begin{equation}
q_\alpha(t + \Delta t) = q_\alpha(t) + p_\alpha\left(t + \Delta t/2\right) \Delta t \,.
\end{equation}
In the second half-step, stochastic forces are added through a modified acceleration:
\begin{equation}
f_\alpha(t + \Delta t) \to f_\alpha(t + \Delta t) - \gamma_\mathrm{c} \, p_\alpha\left(t + \Delta t/2\right) + \xi_\alpha(t) \,,
\end{equation}
where $\gamma_\mathrm{c} = 1/(2\tau_\mathrm{c})$ is the friction coefficient and $\xi_\alpha(t)$ is Gaussian random variable with variance %white noise with variance
\begin{equation}
\langle \xi_\alpha^2 \rangle = \frac{6 \gamma_\mathrm{c} k_\mathrm{B} T}{\Delta t}
\end{equation}
as required by the fluctuation-dissipation theorem.\cite{zwanzig2001nonequilibrium}
The final momentum is then
\begin{equation}
p_\alpha(t + \Delta t) = p_\alpha\left(t + \Delta t/2\right) + \frac{1}{2} f_\alpha(t + \Delta t) \Delta t \,.
\end{equation}
This Langevin integrator maintains the cavity mode at temperature $T$ and thermalizes kinetic energy at a rate $\gamma_\mathrm{c}$. We set the rate as $\gamma_\mathrm{c} = 0.5$ $\mathrm{ps}^{-1}$, ensuring efficient energy exchange with the electromagnetic environment at ultrafast timescales.

For molecular degrees of freedom, we use a two-step velocity-Verlet propagator combined with the stochastic velocity-rescaling algorithm to apply the Bussi-Parrinello bath.\cite{bussi2007canonical,bussi2008stochastic}
At each timestep $\Delta t$, the molecular coordinates $(\bm{R}_i, \bm{P}_i)$ for particle $i$ with mass $M_i$ evolve as follows.
The first half-step advances momenta by
\begin{equation}
\bm{P}_i\left(t + \Delta t/2\right) = \bm{P}_i(t) + \frac{1}{2} \bm{F}_i(t) \Delta t \,,
\end{equation}
where $\bm{F}_i(t)$ is the total force on particle $i$ arising from intermolecular, intramolecular, and cavity coupling interactions.
The momentum is then rescaled by a stochastic factor $\alpha_\upsilon$ computed from the current kinetic energy
\begin{equation}
K_{t+\Delta t/2} = \sum_{i=1}^{N} \frac{|\bm{P}_i(t + \Delta t/2)|^2}{2M_i} \,,
\end{equation}
the temperature $T$, and the thermostat time constant $\tau_\mathrm{b}$, which is set to 1.0 ps, yielding % The scaling factor is
\begin{equation}
\alpha_\upsilon =
\biggl[
e^{-\Delta t/\tau_\mathrm{b}}
+ \frac{T}{2K_{t+\Delta t/2}}
\bigl(1 - e^{-\Delta t/\tau_\mathrm{b}}\bigr)\bigl(R_\Gamma + R_\mathrm{N}^2\bigr)
+ 2R_\mathrm{N}
\sqrt{
\frac{T}{2K_{t+\Delta t/2}}\,
e^{-\Delta t/\tau_\mathrm{b}}\bigl(1 - e^{-\Delta t/\tau_\mathrm{b}}\bigr)
}
\biggr]^{1/2} \,,
\end{equation}
where $R_\mathrm{N}$ is a standard normal random variable, and $R_\Gamma$ is an independent chi-squared random variable with $3N-1$ degrees of freedom.
After rescaling momenta $\bm{P}_i(t + \Delta t/2) \to \alpha_\upsilon \bm{P}_i(t + \Delta t/2)$, positions update according to
\begin{equation}
\bm{R}_i(t + \Delta t) = \bm{R}_i(t) + \frac{\bm{P}_i(t + \Delta t/2)}{M_i} \Delta t \,.
\end{equation}
Forces are then recalculated at the new positions $\bm{F}_i(t + \Delta t)$, after which momenta undergo a second rescaling by a freshly computed factor $\alpha'_\upsilon$ followed by the second half-step:
\begin{equation}
\bm{P}_i(t + \Delta t) = \alpha'_\upsilon \left[\bm{P}_i(t + \Delta t/2) + \frac{1}{2} \bm{F}_i(t + \Delta t) \Delta t\right] \,.
\end{equation}
This integration scheme exactly samples the canonical ensemble while preserving the microscopic dynamics, including ballistic motion at short times, making it the ideal bath for the molecular system. 

\subsection{Adaptive Timestepping}

Simulating systems with disparate timescales, e.g., femtosecond Rabi oscillations to nanosecond structural dynamics, requires careful control of the timestep.
Furthermore, our simulations employ instantaneous step changes in the coupling constant (e.g., a step profile $\lambda(t) = \lambda \Theta(t-t_0)$ or square-wave modulation $\lambda(t) = \lambda\, \mathrm{sq}_D(\Omega_\mathrm{p}(t-t_0))$) that lead to sudden changes in momenta, requiring smaller integration timesteps to stabilize.
To account for these changes in real time, we implement adaptive timestepping, in which the integrator dynamically adjusts $\Delta t$. 
The adaptive timestepping algorithm estimates the local truncation error of the trajectory using a single explicit Euler step as a reference absolute error.

For a particle $i$ at position $\bm{R}_i(t)$ with momentum $\bm{P}_i(t)$ and mass $M_i$, define the velocity $\bm{v}_i(t) = \bm{P}_i(t)/M_i$ and acceleration $\bm{a}_i(t) = \bm{F}_i(t)/M_i$.
The velocity-Verlet predictor yields
\begin{equation}
\bm{R}_{i,\mathrm{VV}}(t + \Delta t) = \bm{R}_i(t) + \bm{v}_i(t) \Delta t + \frac{1}{2} \bm{a}_i(t) (\Delta t)^2 \,,
\end{equation}
while the explicit Euler method gives
\begin{equation}
\bm{R}_{i,\mathrm{E}}(t + \Delta t) = \bm{R}_i(t) + \bm{v}_i(t) \Delta t \,.
\end{equation}
The per-particle error estimate at time $t$ is
\begin{align}
\delta_i(t) &= \left| \bm{R}_{i,\mathrm{VV}}(t + \Delta t) - \bm{R}_{i,\mathrm{E}}(t + \Delta t) \right|  = \frac{1}{2} \left|\bm{a}_i(t)\right| (\Delta t)^2 \,,
\end{align}
and the system-wide error between the two trajectories is computed as the root-mean-square
\begin{equation}
\delta (t) = \sqrt{\frac{1}{N} \sum_{i=1}^{N} \delta_i(t)^2}
\end{equation}
over all $N$ particles. The timestep for time $t$ then follows the formula
\begin{equation}
\Delta t' = \Delta t \sqrt{\frac{\varepsilon^*}{\delta(t)}} \,,
\end{equation}
where $\varepsilon^*$ is the target error tolerance. This approach automatically increases $\Delta t$ when forces are weak and decreases it during high-force events.

\begin{figure}
    \centering
    \includegraphics[width=0.55\linewidth]{coupling.pdf}
    \caption{Schematic of the dynamic error tolerance used in addition to the adaptive timestepping. When the coupling constant exceeds a threshold, we immediately set the error tolerance to nearly zero and ramp it back to its previous value exponentially.}
    \label{fig:placeholder}
\end{figure}
In practice, cavMD simulations reach numerical instabilities due to the step change in our coupling constants, even when implementing a fixed error tolerance. 
To prevent numerical instabilities during protocol switches, we employ error tolerance resetting and ramping. When the coupling constant changes by $\Delta \lambda \geq \varepsilon_\lambda$, we reset the error tolerance to the following exponential ramping function: 
\begin{equation}
\varepsilon(t)= \begin{cases}
\varepsilon^*, & t<t_0, \\ 
\varepsilon^*\left[1-(1-f) e^{-\left(t-t_0\right) / \tau^*}\right], & t \geq t_0 .
\end{cases}
\end{equation}
where $f_0$ is the initial fraction that the error tolerance drops to, and $\tau_*$ is the ramping time constant.
Starting with a conservative timestep at $t=0$, the integrator gradually increases $\Delta t$ as $\varepsilon(t)$ approaches $\varepsilon_*$.
In this work, we use $\varepsilon_* = 5.0$, $f_0 = 10^{-3}$, and $\tau_* = 50\text{ ps}$ for all production runs. 
In practice, we have found that these choices lead to a timestep $\Delta t \sim 1.5$ fs in a stable run, which can drop to $\Delta t \sim 10^{-4}$ fs during the initial drop. 
This choice balances numerical stability during the initial transient period against computational efficiency once the system equilibrates.
When $f_0 = 0$, ramping is disabled and the integrator uses fixed $\varepsilon_*$ throughout, which we employ only for well-equilibrated continuation runs.

\subsection{System Parameters and Force Field}

The system parameters, force field, and associated table are given in the Methods section of the main text.

\subsection{Simulation Protocol and Implementation}

The simulation protocol and implementation details are described in the Methods section of the main text. For further details on the \texttt{cavHOOMD-blue} software package, please refer to the documentation website (\href{https://cav-hoomd.readthedocs.io/en/latest/}{https://cav-hoomd.readthedocs.io/en/latest/}).

\section{Note on Experimental Timescale Requirements and Feasibility}
\label{sec:experimental_feasibility}
% [Muhammad: Done - Expanded with experimental feasibility (cavity/sample, coupling, multimode, structural state)]
%\paragraph*{Timescales.}
While our MD simulations operate the C$^2$F protocol on picosecond timescales, this is chosen only because molecular dynamics simulations are computationally limited to nanosecond ranges. The C$^2$F protocol's only fundamental requirement is \emph{separation of timescales}: cavity and bath modulation must occur much faster than structural relaxation ($\tau_\mathrm{mod} \ll \tau_\mathrm{s}$) and thus we can operate the protocol at lower temperatures. For instance, in deeply supercooled liquids we have $\tau_\mathrm{s} \gtrsim$ milliseconds, and thus microsecond cavity control provides a three-order-of-magnitude separation in timescales, more than sufficient for the protocol to function as designed.
Experimentally, several techniques can achieve these timescales. Electro-optic modulation via Pockels cells enables nanosecond-to-microsecond switching of cavity parameters \cite{dutta2017alloptical,gryaznov2019active}. Tunable-finesse optical resonators based on coupled Fabry-Perot cavities achieve microsecond response times through controlled impedance mismatch \cite{stone2020optical}. For slower glass-forming systems, liquid crystal infiltration into cavities provides millisecond tuning \cite{yavorskiy2025terahertz}, while MEMS and piezoelectric mirrors offer intermediate timescales. Additionally, plasmonic nanoparticle lattices demonstrate tunable light-matter coupling through refractive index modulation \cite{wang2019manipulating}, potentially enabling even faster control.

Separation of timescales also implies that the cavity does not need to be in resonance. As shown in \cref{app:offres} and \cref{fig:fig5}, scanning the cavity frequency $\omega_\mathrm{c}$ around the molecular vibration frequency reveals a weak monotonic frequency dependence of the relaxation slowdown. In other words, non-thermal aging persists at off-resonant frequencies as long as the timescale separation condition is maintained. Additionally, fictive temperature $T_\mathrm{f}$ can be measured experimentally through multiple complementary techniques. Standard differential scanning calorimetry (DSC) and fast scanning calorimetry (FSC) provide direct access to $T_\mathrm{f}$ with measurement timescales ranging from minutes to milliseconds \cite{elbanna2023challenging}. Spectroscopic methods based on fluctuation-dissipation relations offer faster probes; these well-established techniques use the ratio of response-to-fluctuation spectra in the frequency domain to extract the structural fictive temperature,\cite{grigera1999observation,henot2024crossing,zhang2021unifying} and are potentially attractive for incorporating fictive temperature measurements in Fabry-Perot cavities. %These enable experimental validation of the deeply equilibrated states predicted by the C$^2$F protocol.


%\paragraph*{Cavity and sample geometry.}
Our proposal to confine supercooled liquids within optical cavities using silica spacers is experimentally standard. Fused silica, ultra-low-expansion (ULE) glass, and Zerodur are routinely used as cavity spacer materials due to their low thermal expansion and chemical inertness \cite{berthold1977dimensional,legero2010tuning,smid2012evaluation}. Optofluidic Fabry-P\'erot cavities with fused silica capillaries have been demonstrated for liquid samples \cite{gaber2015volume}, and both spacerless flow-cell designs \cite{yamada2022spacerless} and standard cavity configurations \cite{fukushima2022inherent} have successfully achieved vibrational strong coupling in liquid-phase systems including water and ionic liquids. Cavity lengths of order microns to match IR vibrational modes are standard; sample thickness can be matched to the cavity mode profile to maximize coupling.

%\paragraph*{Coupling strength and multimode effects.}
Achievable vacuum Rabi splittings in vibrational strong coupling range from tens to hundreds of cm$^{-1}$, corresponding to coupling strengths comparable to or exceeding those used here.\cite{li2020cavity,fukushima2022inherent} Multimode cavity effects can be minimized by operating with a single dominant mode or by ensuring sufficient frequency separation between cavity modes and molecular vibrations; our simulations retain only the fundamental mode, consistent with this regime.

%\paragraph*{Access to structural state.}
Lastly, for C$^2$F cooling, the feedback loop requires periodic measurement of $T_\mathrm{s}$; the techniques described above (DSC, FSC, and spectroscopic fluctuation-dissipation methods) provide the necessary access, with acquisition times short compared to structural relaxation.

\section{The \texorpdfstring{C$^2$F}{C2F} Protocol Implementation Details}
\label{app:c2f_protocol}

In this section, we explain how we implement the cavity configurational feedback (C$^2$F) cooling protocol. The protocol combines three ingredients: (1) periodic on--off modulation of the cavity coupling $\lambda(t)$, (2) measurement of the structural fictive temperature $T_\mathrm{s}(t)$, and (3) feedback control of the bath temperature $T$ so that it tracks $T_\mathrm{s}$. We describe each in turn, then present the simulation parameters and results.

The structural fictive temperature $T_\mathrm{s}(t) := \vartheta^{-1}[V_\mathrm{LJ}(t)+V_\mathrm{C}(t)]$ is obtained by inverting the equilibrium structural potential-energy calibration $\vartheta(T) := \langle V_\mathrm{LJ}+V_\mathrm{C}\rangle$ based on Rosenfeld--Tarazona scaling,\cite{rosenfeld1998density} as described in the main text. The feedback controller adjusts the bath temperature to follow $T_\mathrm{s}$ by defining the error $e(t) = T(t)-T_\mathrm{s}(t)$ and evolving the bath temperature as
\begin{equation}
\frac{\mathrm{d}T(t)}{\mathrm{d} t} = -k_\mathrm{I}\,e(t) \,,
\label{eq:c2f_feedback}
\end{equation}
where $k_\mathrm{I}$ is the control gain. When $T_\mathrm{s}(t)$ varies slowly, \cref{eq:c2f_feedback} reduces to a PID controller, with only the integral term on, for temperature regulation. The "on-off" effect of the coupling strength is simulated via a square wave:
\begin{equation}
\lambda(t) = \lambda \; \mathrm{sq}_D\!\bigl(\Omega_\mathrm{p}(t-t_\mathrm{0})\bigr) \,,
\label{eq:c2f_squarewave}
\end{equation}
where $\mathrm{sq}_D(\cdot)$ has duty cycle $D$ and period $\Omega_\mathrm{p}^{-1}$. We choose $D=10\,\%$ and $\Omega_\mathrm{p}^{-1}=10$~ps, so the coupling is active for 1~ps per cycle and off for 9~ps. 
%The short on-phase reflects the $O(1)$~ps lifetime of cavity modes observed experimentally.\cite{D4CS01024H,10.1063/5.0282687} 
The only fundamental requirement is $\tau_\mathrm{mod} \ll \tau_\mathrm{s}$; the precise duty cycle and period can be tuned provided this timescale separation holds.

Each cooling cycle proceeds as follows. When the cavity switches on, energy is pumped into the fast vibrational modes, driving the structural degrees of freedom into deeper potential-energy configurations and lowering $T_\mathrm{s}$. The controller then reduces the bath temperature $T$ to match $T_\mathrm{s}$ via \cref{eq:c2f_feedback}. When the cavity switches off, the structural degrees of freedom remain in their deeper configurations because they relax on much longer timescales. The cavity then switches back on from this new, lower structural state. Repeating this cycle progressively brings the system to lower configurational states until a steady state is reached where $T \approx T_\mathrm{s}$, avoiding the kinetic arrest that limits conventional cooling.

For the simulations, initial configurations are first prepared by equilibrating the system at $T = 300$~K without the cavity for 150~ns. From this trajectory, 500 independent configurations separated by 300~ps are extracted as initial conditions for production runs. At $t=0$, a cavity mode is introduced with position and momentum drawn from its thermal equilibrium distribution, and the coupling is first activated at $t_0 = 20$~ps. Both baths operate at temperature $T$ with time constants $\tau_\mathrm{b} = \tau_\mathrm{c} = 1.0$~ps ($\gamma_\mathrm{b}=\gamma_\mathrm{c}=0.5$~ps$^{-1}$). Energy decompositions are recorded every 0.1~ps.

The protocol cools a room-temperature liquid to $T_\mathrm{g} \approx 32$~K within $\sim 50$~ps, producing ultralow-temperature equilibrium configurations inaccessible through conventional cooling. The rapid timescale reflects fast cavity modulation accessible in MD simulations; the underlying mechanism relies only on $\tau_\mathrm{mod} \ll \tau_\mathrm{s}$ and therefore applies equally to slower supercooled liquids at experimentally accessible modulation rates, with the material evolving on correspondingly longer timescales (see \cref{sec:experimental_feasibility}).

\section{Supplementary Results}

\begin{figure}
    \centering
    \includegraphics[width=0.85\linewidth]{combined_fkt_arrhenius.pdf}
    \caption{Equilibrium relaxation behavior of the molecular Kob-Andersen model. (a) The normalized intermediate scattering function at equilibrium, $\phi_{\bm k,\mathrm{eq}}(t)$, plotted over time from $T = 200$ K to $T =65$ K. (b) The equilibrium relaxation time, defined as when $\phi_{\bm k,\mathrm{eq}}(\tau_\mathrm{s}) = 0.1$, along with the Arrhenius curve (red), parabolic law curve (green), the onset temperature (red dashed line, star symbol). Note that the data here is at zero coupling, i.e., with no cavity.}
    \label{fig:feqdata}
\end{figure}
\subsection{Equilibrium Relaxation Behaviors} \label{app:equildata}

The equilibrium behavior of the supercooled liquid establishes a baseline against which aging-induced changes can be measured.
\Cref{fig:feqdata}a shows the intermediate scattering function at equilibrium from $T = 200$~K to $T = 65$~K; relaxation extends over multiple decades in time as the temperature drops by more than 100~K.
When computing the relaxation time, shown in \cref{fig:feqdata}b, we observe the typical behavior: it follows an Arrhenius law at high temperatures, then deviates at lower temperatures. 
Many functional forms can fit the relaxation time; here, we choose the parabolic law. Given an onset temperature $T_\mathrm{o}$ for glassy dynamics, the parabolic law states that\cite{Elmatad2009CorrespondingStates}
\begin{equation}
\ln \left[\frac{\tau_\mathrm{s,eq}(T)}{\tau_\mathrm{o}}\right] = \begin{cases}
(\beta-\beta_\mathrm{o}) E_a  &\beta < \beta_\mathrm{o}
\\
(\beta-\beta_\mathrm{o}) E_a+J^2(\beta-\beta_\mathrm{o})^2  &\beta \geq \beta_\mathrm{o}
\end{cases}
\end{equation}
Within the so-called dynamical facilitation perspective of glassy dynamics, $J$ is proportional to the activation energy $J_\sigma$ of excitations, which are localized pure-shear deformations that mediate structural rearrangements in supercooled liquids.\cite{hasyim2021theory}
At low temperatures, excitations typically facilitate the creation of other excitations nearby. Such a cooperative mechanism leads to non-Arrhenius behavior in the dynamics of density-mode fluctuations.\cite{hasyim2024emergent}

\begin{figure}[t]
    \centering
    \includegraphics[width=0.6\linewidth]{Figure5.pdf}
    \caption{Non-thermal aging at different cavity frequencies. (a) The normalized structural relaxation time over (a) frequency $\omega_\mathrm{c}$ and (b) waiting time $t_\mathrm{w}$, showing the non-monotonic behavior across $\omega_\mathrm{c}$ and $t_\mathrm{w}$.}
    \label{fig:fig5}
\end{figure}


% [Muhammad: Done - Added Kob-Andersen ref]
For the molecular Kob--Andersen model~\cite{kob1995testing}, we obtain from the parabolic-law fitting that the relaxation time at onset $\tau_\mathrm{o} = 24.83$ ps and at a temperature of $T_\mathrm{o} = 134.06$ K, which is slightly higher than the freezing temperatures or triple points of oxygen and nitrogen molecules that the model is based on.\cite{Stewart1991Oxygen,Lemmon2000Air}
The Arrhenius activation energy is $E_\mathrm{a} = 0.049$ eV, which translates to $\sim 4 k_\mathrm{B} T$ at the onset temperature. 
Such an activation energy is moderately high as the system enters its supercooled regime, but practically negligible at room temperature. 
The energy scale $J \sim E_\mathrm{a}$; experimentally, this is a typical result but does not forbid the system to be kinetically arrested at lower temperatures.\cite{Elmatad2009CorrespondingStates}
In fact, the glass transition temperature, defined as when the relaxation time $\tau_\mathrm{s,eq}(T_\mathrm{g}) = 100$ seconds, occurs at $T_\mathrm{g} = 32$ K---only a 100 Kelvin drop from the onset temperature. At such temperatures, it is practically impossible to relax the supercooled liquid to its equilibrium state using brute-force molecular dynamics (MD) simulations.



\subsection{The Role of Polaritons} \label{app:offres}

The mechanism for nonthermal aging relies heavily on strong light-matter interaction and separation of timescales.
Consequently, one might expect that resonance between the cavity mode and intramolecular vibrations---i.e., polariton formation---is not strictly required. Indeed, the resonant condition used in the main text serves primarily to \emph{spectroscopically confirm} strong coupling through the observed Rabi splitting, and we have explained the main text that the mechanism behind non-thermal aging is more general. To see this, we can analyze \cref{eq:HF-classical} and observe that the light-matter interaction shifts the equilibrium position of each cavity coordinate by $q_{0,\alpha} = -\lambda \mu_\alpha / \omega_\mathrm{c}$, a finite displacement whose magnitude depends on $\lambda$ and on the instantaneous dipole moment $\mu_\alpha$, but not on whether $\omega_\mathrm{c}$ is resonant with a vibrational mode. This displacement injects energy into the molecular system through the dipole self-energy term $\lambda^2 \mu_\alpha^2 / (2\omega_\mathrm{c})$, and the resulting energy redistribution drives the structural degrees of freedom to deeper potential minima regardless of resonance. Resonant conditions are therefore not required for non-thermal aging; what matters is the energy added via the finite-$q$ shift and the separation of timescales between fast light-matter dynamics and slow structural relaxation. Nevertheless, the subsequent dynamics still depend on $\omega_\mathrm{c}$ and $\lambda$ in a non-trivial way.

To test the importance of resonance conditions, we fix the coupling strength at $\lambda = 0.141$~a.u.\ and scan cavity frequencies around the molecular vibration frequency $\omega_\mathrm{AA} = 1560~\mathrm{cm}^{-1}$. 
If polariton formation were essential, we would expect a sharp peak in the normalized structural relaxation time $\tilde{\tau}_\mathrm{s}$ at $\omega_\mathrm{c} = \omega_{\mathrm{AA}}$, followed by rapid decay to the equilibrium baseline ($\tilde{\tau}_\mathrm{s}=1$) at off-resonant frequencies.
However, results in \cref{fig:fig5}a reveal a monotonic effect where the initial slowdown is more pronounced at low frequencies before tapering off at high frequencies.
Furthermore, as shown in \cref{fig:fig5}b, aging at low frequencies retains memory longer, reaching equilibrium at longer waiting times than higher frequencies. Notably, some conditions display non-monotonic behavior in waiting time, with maximum slowdown at intermediate $t_\mathrm{w} = 200$~ps.

These results confirm that non-thermal aging does not require polariton formation.
The mechanism depends on strong light--matter coupling in general, with the magnitude and temporal characteristics of aging depending non-trivially on $\omega_\mathrm{c}$ and $\lambda$.
All results here use the intermediate scattering function at a fixed wavevector $|\bm{k}| = 6.02$~a.u.; cavity-mediated effects could in principle introduce a finite $q$-dependent shift, but our data indicate that the qualitative aging behavior is captured at this wavevector.


\subsection{Benchmarking \texttt{cavHOOMD-blue} against LAMMPS+i-PI}

\begin{figure}[t]
    \centering
    \includegraphics[width=0.85\linewidth]{combined_analysis.pdf}
    \caption{Validation of \texttt{cavHOOMD-blue} against LAMMPS+i-PI~\cite{li2020cavity}. 
    (a) Normalized dipole autocorrelation function and (b) infrared absorption spectrum for the molecular Kob-Andersen model at equilibrium ($T = 100$~K, $\omega_\mathrm{c} = 1560$~cm$^{-1}$). 
    Dashed vertical lines indicate the bare vibrational frequencies: A$_2$ stretch at 1560~cm$^{-1}$ and B$_2$ stretch at 2335~cm$^{-1}$.}
    \label{fig:cavmd_benchmark}
\end{figure}

To validate our GPU-accelerated implementation of cavity molecular dynamics in \texttt{cavHOOMD-blue}, we benchmark against the established LAMMPS+i-PI framework developed by Li, Subotnik, and Nitzan~\cite{li2020cavity}. The original cavity molecular dynamics (CavMD) approach was introduced in Ref.~\cite{li2020cavity}, where the authors demonstrated vibrational strong and ultrastrong coupling in liquid water, revealing asymmetric Rabi splitting and cavity-modified orientational dynamics. Their implementation couples the i-PI universal force engine~\cite{ipi2014} to LAMMPS for the molecular dynamics, providing a reference standard for CavMD simulations. We compare the equilibrium properties computed by both codes using identical simulation parameters: the molecular Kob-Andersen model at $T = 100$~K with cavity frequency $\omega_\mathrm{c} = 1560$~cm$^{-1}$ resonant with the A$_2$ stretch.

\Cref{fig:cavmd_benchmark} shows the comparison of the normalized dipole autocorrelation function $C_{\mu\mu}(t)/C_{\mu\mu}(0)$ and the infrared absorption spectrum $n(\omega)\alpha(\omega)$ computed by \texttt{cavHOOMD-blue} and LAMMPS+i-PI. Both quantities exhibit excellent agreement between the two implementations. The dipole autocorrelation functions overlap throughout the decay, confirming that the cavity-molecule coupling and thermostatting schemes produce identical equilibrium dynamics. The infrared spectra display the characteristic peaks at the A$_2$ stretch (1560~cm$^{-1}$) and B$_2$ stretch (2335~cm$^{-1}$) frequencies, with matching lineshapes and intensities. This quantitative agreement confirms that \texttt{cavHOOMD-blue} reproduces the original CavMD framework while providing the GPU acceleration needed for the long-timescale simulations of aging studied in this work.
